% GENERATED FILE. Edit canonical lessons, then run npm run generate:books.
% canonical-source-hash: fnv1a64:e540f13c0d88371b
% canonical-lessons: RU-C473-razocharovaniye, RU-C473-razocharovannyy, RU-C473-oblegcheniye, RU-C473-otchayaniye, RU-C473-vzdykhat

\chapter{Disappointment and Relief}
\label{ch:ru-disappointment-and-relief}
\hlchaptermodality{473}

\begin{chapteropening}
\textbf{By the end of this chapter:} \emph{I can talk about disappointment, despair and relief.}

Complete the last lesson of chapter 473: I can talk about disappointment, despair and relief.
\end{chapteropening}

\section[\texorpdfstring{\ru{разочарование} (razocharovániye)}{разочарование (razocharovániye)}]{\ru{разочарование} (razocharovániye) — disappointment}

\label{lesson:RU-C473-razocharovaniye}



\begin{quote}
Before the new one: say the Russian for fury, then the Russian for quick-tempered.
\end{quote}

\subsection*{You'll want to know: \ru{разочарование}}
\textbf{\ru{разочарование}} (\emph{razocharovániye}) — \textquotedblleft{}disappointment\textquotedblright{}.

\begin{quote}
\ru{Фильм} \ru{стал} \ru{для} \ru{нас} \ru{большим} \ru{разочарованием}.
\end{quote}

\begin{quote}
\emph{Film stal dlya nas bolshím razocharovániyem.}
\end{quote}

\begin{quote}
The film was a big disappointment for us.
\end{quote}

Neuter. The verbs are \emph{\ru{разочаровываться} / \ru{разочароваться}} — to be disappointed.

\begin{cousinweb}
\emph{\ru{разочарование}} is built on \emph{\ru{разочаровать}} 'to break the spell', from \emph{\ru{очаровать}} 'to enchant' and \emph{\ru{чары}} 'spells'; EtymDB traces \emph{\ru{чары}} to Proto-Slavic \emph{čar\ru{ъ}} 'magic, sorcery'.
\end{cousinweb}

\begin{grammarlens}[title={What you've built}]
One more word for telling, explaining and describing at B1.
\end{grammarlens}

\subsection*{Your turn}
\begin{itemize}
\raggedright
  \item \emph{Say it:} \emph{razocharovániye}
  \item \emph{Say it:} \emph{razocharovániye}, once more
  \item \emph{Say it:} a sentence of your own with \emph{razocharovániye}
  \item \emph{Recall:} say the Russian for fury, then the Russian for quick-tempered, then say \emph{razocharovániye} again
\end{itemize}

\begin{tcolorbox}[breakable,colback=teal!4,colframe=teal!35!black,title={Before you move on}]
What does \ru{разочарование} mean? (\textquotedblleft{}Disappointment\textquotedblright{}.) Say it once more.
\end{tcolorbox}
\section[\texorpdfstring{\ru{разочарованный} (razocharóvannyy)}{разочарованный (razocharóvannyy)}]{\ru{разочарованный} (razocharóvannyy) — disappointed}

\label{lesson:RU-C473-razocharovannyy}



\begin{quote}
Before the new one: say the Russian for quick-tempered, then the Russian for disappointment.
\end{quote}

\subsection*{You'll want to know: \ru{разочарованный}}
\textbf{\ru{разочарованный}} (\emph{razocharóvannyy}) — \textquotedblleft{}disappointed\textquotedblright{}.

\begin{quote}
\ru{Он} \ru{вышел} \ru{из} \ru{кабинета} \ru{разочарованный}.
\end{quote}

\begin{quote}
\emph{On výshel iz kabinéta razocharóvannyy.}
\end{quote}

\begin{quote}
He came out of the office disappointed.
\end{quote}

Past participle of \emph{\ru{разочаровать}}, used as an adjective. Short form: \emph{\ru{разочарован} \ru{в} \ru{ком}} — disappointed in someone.

\begin{cousinweb}
\emph{\ru{разочарованный}} is the past participle of \emph{\ru{разочаровать}}, built on \emph{\ru{очаровать}} 'to enchant' and \emph{\ru{чары}} 'spells', which EtymDB traces to Proto-Slavic \emph{čar\ru{ъ}} 'magic'.
\end{cousinweb}

\begin{grammarlens}[title={What you've built}]
One more word for telling, explaining and describing at B1.
\end{grammarlens}

\subsection*{Your turn}
\begin{itemize}
\raggedright
  \item \emph{Say it:} \emph{razocharóvannyy}
  \item \emph{Say it:} \emph{razocharóvannyy}, once more
  \item \emph{Say it:} a sentence of your own with \emph{razocharóvannyy}
  \item \emph{Recall:} say the Russian for quick-tempered, then the Russian for disappointment, then say \emph{razocharóvannyy} again
\end{itemize}

\begin{tcolorbox}[breakable,colback=teal!4,colframe=teal!35!black,title={Before you move on}]
What does \ru{разочарованный} mean? (\textquotedblleft{}Disappointed\textquotedblright{}.) Say it once more.
\end{tcolorbox}
\section[\texorpdfstring{\ru{облегчение} (oblegchéniye)}{облегчение (oblegchéniye)}]{\ru{облегчение} (oblegchéniye) — relief}

\label{lesson:RU-C473-oblegcheniye}



\begin{quote}
Before the new one: say the Russian for disappointment, then the Russian for disappointed.
\end{quote}

\subsection*{You'll want to know: \ru{облегчение}}
\textbf{\ru{облегчение}} (\emph{oblegchéniye}) — \textquotedblleft{}relief\textquotedblright{}.

\begin{quote}
\ru{Когда} \ru{врач} \ru{сказал}, \ru{что} \ru{всё} \ru{в} \ru{порядке}, \ru{я} \ru{почувствовал} \ru{облегчение}.
\end{quote}

\begin{quote}
\emph{Kogdá vrach skazál, chto vsyo v poryádke, ya pochúvstvoval oblegchéniye.}
\end{quote}

\begin{quote}
When the doctor said everything was fine, I felt relief.
\end{quote}

Neuter. \emph{\ru{С} \ru{облегчением}} — with relief.

\begin{cousinweb}
\emph{\ru{облегчение}} is built on \emph{\ru{облегчить}} 'to lighten', from \emph{\ru{лёгкий}} 'light, easy'; EtymDB lists \emph{\ru{облегчить}} without a source, and \emph{\ru{лёгкий}} is inherited from Proto-Slavic \emph{l\ru{ь}g\ru{ъ}k\ru{ъ}} (EtymDB).
\end{cousinweb}

\begin{grammarlens}[title={What you've built}]
One more word for telling, explaining and describing at B1.
\end{grammarlens}

\subsection*{Your turn}
\begin{itemize}
\raggedright
  \item \emph{Say it:} \emph{oblegchéniye}
  \item \emph{Say it:} \emph{oblegchéniye}, once more
  \item \emph{Say it:} a sentence of your own with \emph{oblegchéniye}
  \item \emph{Recall:} say the Russian for disappointment, then the Russian for disappointed, then say \emph{oblegchéniye} again
\end{itemize}

\begin{tcolorbox}[breakable,colback=teal!4,colframe=teal!35!black,title={Before you move on}]
What does \ru{облегчение} mean? (\textquotedblleft{}Relief\textquotedblright{}.) Say it once more.
\end{tcolorbox}
\section[\texorpdfstring{\ru{отчаяние} (otcháyaniye)}{отчаяние (otcháyaniye)}]{\ru{отчаяние} (otcháyaniye) — despair}

\label{lesson:RU-C473-otchayaniye}



\begin{quote}
Before the new one: say the Russian for disappointed, then the Russian for relief.
\end{quote}

\subsection*{You'll want to know: \ru{отчаяние}}
\textbf{\ru{отчаяние}} (\emph{otcháyaniye}) — \textquotedblleft{}despair\textquotedblright{}.

\begin{quote}
\ru{В} \ru{отчаянии} \ru{она} \ru{позвонила} \ru{сестре}.
\end{quote}

\begin{quote}
\emph{V otcháyanii oná pozvoníla sestré.}
\end{quote}

\begin{quote}
In despair she phoned her sister.
\end{quote}

Neuter. The verbs are \emph{\ru{отчаиваться} / \ru{отчаяться}} — to despair.

\begin{cousinweb}
\emph{\ru{отчаяние}} is built on \emph{\ru{отчаяться}} 'to give up hope', from \emph{\ru{от}-} and \emph{\ru{чаять}} 'to hope', which is inherited from Proto-Slavic \emph{čajati} 'to expect, to wait' (EtymDB).
\end{cousinweb}

\begin{grammarlens}[title={What you've built}]
One more word for telling, explaining and describing at B1.
\end{grammarlens}

\subsection*{Your turn}
\begin{itemize}
\raggedright
  \item \emph{Say it:} \emph{otcháyaniye}
  \item \emph{Say it:} \emph{otcháyaniye}, once more
  \item \emph{Say it:} a sentence of your own with \emph{otcháyaniye}
  \item \emph{Recall:} say the Russian for disappointed, then the Russian for relief, then say \emph{otcháyaniye} again
\end{itemize}

\begin{tcolorbox}[breakable,colback=teal!4,colframe=teal!35!black,title={Before you move on}]
What does \ru{отчаяние} mean? (\textquotedblleft{}Despair\textquotedblright{}.) Say it once more.
\end{tcolorbox}
\section[\texorpdfstring{\ru{вздыхать} (vzdykhát)}{вздыхать (vzdykhát)}]{\ru{вздыхать} (vzdykhát) — to sigh}

\label{lesson:RU-C473-vzdykhat}



\begin{quote}
Before the new one: say the Russian for relief, then the Russian for despair.
\end{quote}

\subsection*{You'll want to know: \ru{вздыхать}}
\textbf{\ru{вздыхать}} (\emph{vzdykhát}) — \textquotedblleft{}to sigh\textquotedblright{}.

\begin{quote}
\ru{Бабушка} \ru{вздыхает} \ru{и} \ru{смотрит} \ru{в} \ru{окно}.
\end{quote}

\begin{quote}
\emph{Bábushka vzdykháyet i smótrit v oknó.}
\end{quote}

\begin{quote}
Grandma sighs and gazes out of the window.
\end{quote}

Imperfective; perfective \ru{вздохнуть}. The noun is \emph{\ru{вздох}} — a sigh.

\begin{cousinweb}
\emph{\ru{вздыхать}} is built on \emph{\ru{вз}-} 'up' and the root of \emph{\ru{дышать}} 'to breathe' and \emph{\ru{дух}} 'breath', which EtymDB traces to Proto-Slavic \emph{dux\ru{ъ}} 'breath, spirit'.
\end{cousinweb}

\begin{grammarlens}[title={What you've built}]
One more word for telling, explaining and describing at B1.
\end{grammarlens}

\subsection*{Your turn}
\begin{itemize}
\raggedright
  \item \emph{Say it:} \emph{vzdykhát}
  \item \emph{Say it:} \emph{vzdykhát}, once more
  \item \emph{Say it:} a sentence of your own with \emph{vzdykhát}
  \item \emph{Recall:} say the Russian for relief, then the Russian for despair, then say \emph{vzdykhát} again
\end{itemize}

\begin{tcolorbox}[breakable,colback=teal!4,colframe=teal!35!black,title={Before you move on}]
What does \ru{вздыхать} mean? (\textquotedblleft{}To sigh\textquotedblright{}.) Say it once more.
\end{tcolorbox}
